% 06-experiences-resultats (en): every experiment's result, ten seeds first, then twenty seeds and equal budget.
\chapter{The experiments: chance fools the eye, sorting fools the test}\label{chap:06-experiences-resultats}\chaptermark{The experiments: results}

Chance in the random seeds fools the eye, rarely the test. Sorting the runs fools the test. Keeping the best
three runs out of ten inflates the mean by a third. A test that counts episodes instead of random seeds
recovers the t of 9 published in 2018. The leading program depends on the task, the budget and the physics.
Five random seeds read a small difference backwards nearly one time in three. A published tuning recipe has
aged with its libraries. This chapter gives the result of each experiment of \cref{chap:05-experiences-methode},
in the same order: the question, the plan in one sentence, the result with its intervals, then what the
result means and what it leaves open.

\paragraph{Reading conventions}
The main results are those with ten random seeds (0 to 9), as the protocol announced. The results with twenty
random seeds (0 to 19) and at equal budget come next. Every divergence is stated. A score is the mean return
of the ten final evaluation episodes of a run (\cref{sec:evaluation}). The score of a program is an
interquartile mean (IQM) of these scores over the random seeds. An interval written in square brackets, such
as $[\num{2268}, \num{3369}]$, is a 95\% bootstrap interval, unless the text names another level. A difference
is \emph{settled} when it satisfies the settled-winner rule (\cref{sec:regle}): the IQM intervals are disjoint,
or the interval of the probability of improvement excludes 0.5. The generated tables use the same word. The
rule exceeds its nominal 5\%. It errs slightly more than one time in twenty at ten runs against ten, nearly
one time in ten at five against five (a control written after the results, \cref{tab:regle-tranche}). The condition on
the probability of improvement produces almost all of these false positives. A settled difference whose interval
comes close to 0.5 is read with this margin in mind. The tables and the numbers in the text are read from
the repository's result files (\cref{chap:08-logiciel}).

\section{E1: chance rarely fools the test, sorting does}\label{sec:e1}

\paragraph{Question and plan}
Can two groups of runs of the same program, which differ only in their random seeds, pass for two different
algorithms? E1 trains PPO PyTorch thirty times on \task{HalfCheetah-v5}, one million steps each, random seeds
0 to 29, then compares groups drawn from these thirty runs with Welch's test (\cref{sec:plan-e1}).

\paragraph{A wide spread that is not bell-shaped}
The mean of the thirty scores is \num{1688}, the standard deviation 568, or 34\% of the mean. The median is
\num{1580}. Twenty-seven runs end between \num{1138} and \num{1999}. Three take off beyond \num{3000}, up to
\num{3430} (\cref{fig:e1-graines}). These three runs pull the mean above the median. The Shapiro-Wilk test
clearly rejects the normal distribution ($p = 4\times10^{-7}$). A t-test assumes roughly normal group means.
With five runs drawn from such a distribution, the assumption is fragile.

\fig{e1-graines}{Thirty trainings of the same program on the same task: PPO written in PyTorch
(Stable-Baselines3) on \task{HalfCheetah-v5}, one million steps each. Each dot is a run, placed at its final
evaluation score, the mean return of ten episodes played without randomness on random seeds never seen in
training. Only the training random seed changes from one dot to the next. Mean \num{1688}, median
\num{1580} (dotted line). Twenty-seven runs end between \num{1138} and \num{1999}, three beyond \num{3000}.}

\paragraph{The test's false positives}
\Cref{tab:e1} and \cref{fig:e1-faux-positifs} give the share of the \num{10000} draws in which Welch's test,
at the 5\% level, declares two groups from the same program different. Each of these declarations is a false
positive.
\begin{itemize}
  \item Three runs against three, drawn at random: 1.6\% of the draws.
  \item Five against five: 1.2\%.
  \item Ten against ten: 1.7\%.
  \item The five best against the five worst of a draw of ten runs: 31.0\%.
\end{itemize}
Drawn at random, the groups fool the test less often than the nominal 5\% level. The test is conservative on
this skewed distribution. The most likely explanation: an outlying run inflates the variance of the group
that receives it along with its mean. Sorted, the groups fool the test nearly one time in three on this program-task pair.
The median $p$-value of the sorted split is 0.09.

\gentable{e1}

\fig{e1-faux-positifs}{Share of the \num{10000} draws in which Welch's test, at the 5\% level, declares two
groups of runs of the same program different (PPO PyTorch on \task{HalfCheetah}, thirty runs that differ only
in their random seed). Each grid has 100 cells, one cell standing for one hundredth of the draws. The cells
framed with a dotted line mark the 5\% level. Drawn at random, three against three, five against five or ten
against ten, the groups fool the test less often than this level. The five best set against the five worst of
a draw of ten fool it nearly one time in three.}

\paragraph{Chance fools the eye}
The groups drawn at random rarely fool the test, yet their means differ a great deal. The relative difference
between the means, divided by the smaller of the two, has a median of 17\% for three against three, 17\% for
five against five and 11\% for ten against ten. One time in twenty (the 95th percentile), it exceeds 62\%, 47\%
and 33\%. Two groups of five runs drawn at random therefore often show well-separated final scores, although
they come from the same program. In the sorted split, the median difference reaches 42\%.

\paragraph{Keeping the best three runs out of ten}
Over \num{10000} draws of ten runs, the mean of the best three averages \num{2249}, against \num{1692} for all
ten: an inflation of 33\%. Experiment E4 extends this measurement to the four programs (\cref{sec:e4}).

\paragraph{The 2018 statistic, out of PPO's reach}
\textcite{henderson2018matters} test whole learning curves. The second statistic of E1, added after the
previous results, recomputes their average t on our thirty curves: a Welch test at each of the
twenty evaluation points, then the mean of the $t$ values (\cref{sec:t-moyen}). One time in twenty, the
absolute value of the average t exceeds 1.4 when the groups of five are drawn at random, 2.5 when they are
sorted by final score. The \num{20000} draws all stay below 9, the value published in 2018 that \cref{sec:b1} examines. The test at the last point alone recovers the
result on final scores: about 31\% false positives for sorted groups.

\paragraph{Sixteen pairs instead of one}
E1 covers one program and one task, whereas the shape of the distribution depends on both. An analysis
decided after the results and written before the computation repeats the two five-against-five draws on each
of the sixteen (program, task) pairs of E2, with their twenty random seeds (\cref{tab:e1-grille},
\cref{fig:e1-grille}).
\begin{itemize}
  \item Drawn at random, the groups fool the test in 2.9\% of the draws at the median, from 0.7\% to 4.8\%
    depending on the pair. The rate stays below the nominal 5\% on all sixteen pairs. It is lowest on the
    most skewed pairs: 0.7\% for PPO JAX on \task{HalfCheetah} (skewness 1.66), 1.0\% for PPO PyTorch on the
    same task (2.08). On nearly symmetric distributions it comes closer to 5\%, as the skewness explanation led
    us to expect.
  \item Sorted, the five best against the five worst fool the test in 83\% of the draws at the median, from
    26\% to 99.8\% depending on the pair. The preregistered pair, PPO PyTorch on \task{HalfCheetah}, gives the
    lowest rate of the sixteen. On \task{Hopper}, both PPOs exceed 98\%.
\end{itemize}
Recomputed with the grid's generator on its thirty runs, the preregistered pair gives 1.3\% at random and
31.1\% sorted, against 1.2\% and 31.0\% in \cref{tab:e1}: the draws alone account for the difference. The
preregistered result remains that of \cref{tab:e1}.

\fig{e1-grille}{E1 repeated on the sixteen (program, task) pairs, twenty random seeds each. For each pair, the
share of the \num{10000} draws in which Welch's test, at the 5\% level, declares two groups of five runs of the
same program different: groups drawn at random (hollow circle) and the five best against the five worst of
a draw of ten (filled dot). Sorting the runs makes the test conclude wrongly from 26\% to nearly 100\% of the
time depending on the pair, while groups drawn at random stay below 5\%. The preregistered case, PPO PyTorch
on \task{HalfCheetah}, is the least fooled of the sixteen.}

\gentable{e1-grille}

\paragraph{What E1 establishes and its limits}
On these thirty runs, chance in the random seeds fools the eye far more often than the test. A Welch test on
groups drawn at random is cautious, even with three runs per group. The grid of sixteen pairs confirms it.
Sorting the runs turns this chance into false positives: nearly one time in three on the preregistered pair,
more than four times in five on nine of the fifteen others. Selecting the best runs inflates the mean by a
third. The grid is an exploratory analysis, decided after the results.

\section{B1: a test that counts episodes recovers the t of 9}\label{sec:b1}

\paragraph{Question}
Figure~5 of \textcite{henderson2018matters} sets two groups of five TRPO runs against each other on
\task{HalfCheetah}, with the same code and the same settings. It reports $t=-9.09$ and $p=0.0016$. Pineau
stated on stage at ICLR 2018 that the groups are not sorted: \enquote{They're not top five. They're actually just five runs.} (around 28~min~40~s into the video, \cite{pineau2018iclr}; flaw F11, \cref{sec:failles}). Yet E1 shows that with PPO, groups drawn at random never come
near an average t of 9. Where does this figure come from? Three leads remained open: the spread specific to
TRPO, another reading of their caption, or another computation (\cref{sec:plan-b1}).

\paragraph{First lead: TRPO spreads more, without nearing 9}
The thirty TRPO runs, trained with the default values of sb3-contrib, spread far more than the PPO runs. Their
final mean is \num{1951}, their standard deviation 930, or 48\% of the mean, against 34\% for PPO. Even so,
with groups of five drawn at random, the average t over the curve exceeds 1.4 only one time in twenty, as for
PPO. With sorted groups, it exceeds 1.9 one time in twenty. The rule fixed before any TRPO training rules out
the spread of the algorithm as the explanation.

\paragraph{Second lead: stacked points, a faulty computation that still falls short}
This lead takes the third of the readings of the average t (\cref{sec:t-moyen}), the one that stacks the
hundred points of the five curves of each group as if they were independent observations. One time in twenty, its $|t|$ exceeds 5.0 with groups drawn at random, 6.3 with sorted groups.
No draw reaches 9. This reading has a worse defect than its failure. It declares a difference in 53\% of the
random draws with TRPO and in 23\% with PPO, although both groups run the same code. The twenty points of a
curve are not independent: a run that starts well often stays ahead. An analysis written after the results and
before this computation puts a number on this dependence. The lag-1 autocorrelation of the difference between a curve and
the mean curve has a median of $r = 0.71$ for PPO and $0.92$ for TRPO. If this difference follows a first-order
autoregressive process with coefficient $r$, a series of $n$ points is worth about $n(1-r)/(1+r)$ independent
points. A twenty-point curve is then worth about 3.4 for PPO and 0.9 for TRPO.

\paragraph{The published p rules out two readings}
A two-sided Welch test that yields $t=9.09$ and $p=0.0016$ has about 3.4~degrees of freedom. This pair rules
out the stacked points, which have about 198~degrees of freedom and would give a $p$ of the order of
$10^{-16}$. It also rules out the average t converted into a $p$: two groups of five give at least 4~degrees
of freedom (\cref{sec:welch}). The pair rules out these two readings without pointing to the script. That
leaves the mean of the $p$-values, which the script also reports. Such a mean has no defined number of degrees
of freedom: the mean $t$ and the mean $p$ are no longer linked by a Student distribution, so a wide range of
$(t, p)$ pairs is compatible with it. The published pair neither contradicts nor confirms this reading.

\paragraph{Third lead: the 2017 script}
In June 2017, Peter Henderson published a script with a telling name, \file{from\_same\_dist.py}
\parencite{henderson2017samedist}. \Cref{sec:script-henderson} gives its exact computation. In short:
\begin{enumerate}
  \item it averages the curves of the five random seeds of each group, point by point;
  \item at each point, it compares the two mean curves with a Welch test whose variance is the variance
    \emph{of the episodes} of one iteration, divided by the number of episodes;
  \item it reports the mean of the $t$ values and the mean of the $p$ values.
\end{enumerate}
The variability between random seeds, the one the figure was meant to measure, never enters the test. The
script counts episodes instead of independent training runs.

\paragraph{The script finds the t of 9 on our TRPO runs}
Applied to our thirty TRPO curves with $n=10$ episodes, following a plan written before the computation,
the script gives the results of \cref{tab:b1}:
\begin{itemize}
  \item With groups drawn at random, the absolute value of the average t reaches 9.09 in 38\% of the draws.
    Its median is 7.1. The test declares a difference in 43\% of the draws.
  \item With sorted groups, 70\% of the draws reach 9.09.
  \item With $n=5$, the sensitivity analysis, 23\% of the random draws reach 9.09. The test then declares a
    difference in 25\% of the draws.
\end{itemize}
Applied to our thirty PPO curves, the same computation never reaches 9.09. One time in twenty, its average
$|t|$ exceeds 3.6 with random groups. \Cref{fig:t-de-neuf} shows this contrast, \cref{fig:b1-distribution}
the full distributions.

\gentable{b1}

\paragraph{Why TRPO alone gets there}
The denominator of the script is the standard deviation between episodes. The correct denominator is the
standard deviation between random seeds. Write $S$ for the first, $\sigma$ for the second and
$\rho = S/\sigma$ for their ratio, measured at the last evaluation point. A Welch test on five random seeds
per group divides the difference between the means by $\sigma\sqrt{2/5}$. The script divides it by $S\sqrt{2/n}$,
with $n$ episodes. The ratio of the two $t$ values is therefore
\begin{equation}
  \frac{t_{\text{script}}}{t_{\text{seeds}}} = \frac{\sigma\sqrt{2/5}}{S\sqrt{2/n}} = \frac{\sqrt{n/5}}{\rho},
\end{equation}
that is $\sqrt2/\rho$ for our ten episodes and $1/\rho$ for the sensitivity analysis with five. A control
written before this computation checks this identity on each of the \num{10000} draws, to the last digit
(largest relative difference $7\times10^{-16}$). For our TRPO, the standard deviation between the
episodes of one agent is only 7\% of the standard deviation between agents. For our PPO, it is 30\%. The median
factor per draw is 18.1 for TRPO and 4.0 for PPO with ten episodes, 12.8 and 2.8 with five. A TRPO agent scores
almost the same at every episode, while two TRPO agents differ a great deal. In this case, counting episodes
instead of random seeds inflates $t$ by this factor.

\fig{t-de-neuf}{Two groups of five runs drawn at random from thirty, \num{10000} times, for each algorithm and
each way of computing the average $t$ along the learning curve. Each grid has 100 cells, one cell standing for
one hundredth of the draws. It colours the share of draws in which the average $t$, sign ignored, reaches
9.09, the value Henderson and his coauthors published in 2018 for TRPO on \task{HalfCheetah}. Only the script
they published in 2017, which tests the spread between episodes instead of the spread between random seeds,
gets there, applied to TRPO: 38\% of the draws. PPO with this script, PPO and TRPO with our test on random
seeds, TRPO with stacked points: no draw at all.}

\fig{b1-distribution}{Distribution of the absolute value of the average $t$ over \num{10000} draws of two
groups of five runs taken at random from thirty, on \task{HalfCheetah}, for TRPO (top) and for PPO PyTorch
(bottom), on a logarithmic scale. In grey, the Welch test computed between random seeds at each point of the
curve, then averaged. In red, the computation of Henderson's script (2017), which averages the curves of each
group and tests the spread between the ten evaluation episodes. The vertical line marks 9.09, the value
published in 2018. Only the script applied to TRPO exceeds it, in 38\% of the draws.}

\paragraph{What remains uncertain}
We cannot prove that this script produced figure~5. Its name, its output (\enquote{t=\dots, p=\dots}), the
order of magnitude it recovers on our TRPO runs and the published $(t,p)$ pair, which rules out the two other
readings, make it likely. Our conditions differ from theirs, as we stated before the computation. Our
evaluation episodes are played by a deterministic policy. Their progress files record the training episodes
of a stochastic policy, which probably vary more. We have twenty evaluation points instead of their
iterations. Our TRPO comes from sb3-contrib, while their code was different. The order of magnitude comes back
despite these differences.

\paragraph{The t of 9, made by the episodes}
The t of 9 of 2018 needs no sorting of the runs to appear. A test that counts episodes instead of random seeds
produces it on our own TRPO runs. Flaw F11, same code, another random seed, goes deeper than the figure
showed. If this script produced the figure, chance in the random seeds fooled the very test that was meant to
measure it. \Cref{fig:e1-henderson}
brings E1 and B1 together against the published figure.

B1 also exposes a limit of the survey. Criterion X2 counts any error bar as yes, whatever spread the bar
measures. A bar computed over the evaluation episodes of a single agent satisfies X2 while it repeats the
script's computation. The exploratory criterion X8, which asks whether the uncertainty is computed over
training random seeds, shows no detectable rise from 2018 to 2026, while X2 rises clearly
(\cref{sec:sr-principale}).

\fig{e1-henderson}{Our thirty PPO runs and our thirty TRPO runs on \task{HalfCheetah}, compared with the
figure published in 2018. For each way of splitting ten runs into two groups of five, at random or sorted by
final score, the bar gives the absolute value of the average $t$ along the learning curve exceeded one time
in twenty, over \num{10000} draws. The stacked-points bar treats the hundred curve points of each group as
independent observations. The bar for Henderson's script applies his 2017 computation, which tests the spread
between episodes, to our TRPO runs. The bar for Henderson and his coauthors is the value they publish on their
own runs.}

\section{E2: the leading program changes from task to task}\label{sec:e2}

\paragraph{Question and plan}
\textcite{henderson2018matters} showed that the algorithm leading on one task stopped leading on another. E2
redoes the comparison with four programs of 2026 (PPO PyTorch, PPO JAX, SAC, CrossQ) on the four tasks, ten
random seeds each, at each program's budget: one million steps for PPO, \num{200000} for SAC and CrossQ
(\cref{sec:plan-e2}).

\paragraph{Ten random seeds: four tasks, three leading programs}
\Cref{tab:e2} gives the IQM of each program on each task with its interval. \Cref{fig:classement-taches}
divides each IQM by the best one of its task.
\begin{itemize}
  \item \task{HalfCheetah}: SAC leads with \num{5191} $[\num{4431}, \num{5601}]$, ahead of CrossQ, \num{4801}
    $[\num{2427}, \num{7815}]$. The two PPOs follow far behind, \num{1407} for PPO PyTorch and \num{1254} for
    PPO JAX.
  \item \task{Hopper}: PPO JAX leads with \num{3064}, ahead of PPO PyTorch (\num{2877}), CrossQ (\num{2302})
    and SAC (\num{1308}).
  \item \task{Walker2d}: CrossQ leads with \num{5414} $[\num{4994}, \num{5627}]$, far ahead of PPO PyTorch
    (\num{2158}), PPO JAX (\num{1272}) and SAC (\num{1236}).
  \item \task{Swimmer}: PPO JAX leads with 112, ahead of CrossQ (110) and PPO PyTorch (107). SAC stays at 47.
\end{itemize}

\fig{classement-taches}{Score of each program divided by the best score of its task, over ten random seeds:
the bar at 100 marks the leading program. The score of a program is the interquartile mean (the mean of the
middle six scores out of ten) of its final evaluation scores. PPO trains for one million steps, SAC and CrossQ
for \num{200000}. The leading program changes: SAC on \task{HalfCheetah}, PPO written in JAX on \task{Hopper}
and \task{Swimmer}, CrossQ on \task{Walker2d}.}

\paragraph{A settled change, a lead seldom settled}
The leading program differs from task to task. The change is settled under the protocol's rule: SAC, leading
on \task{HalfCheetah}, is beaten in a settled way on \task{Hopper} by both PPOs, on \task{Walker2d} by CrossQ
and PPO PyTorch, on \task{Swimmer} by the other three. The preregistered count finds 14~settled changes. The leading
program, however, beats all the others in a settled way only on \task{Walker2d}. On the other three tasks, ten
random seeds do not separate it from the runner-up: SAC from CrossQ on \task{HalfCheetah}, PPO JAX from PPO
PyTorch on \task{Hopper}, PPO JAX from CrossQ on \task{Swimmer}. The interval of CrossQ on
\task{HalfCheetah}, from \num{2427} to \num{7815}, shows how much its runs vary on this task.

\paragraph{With Bonferroni's correction}
E2 computes 24~probability-of-improvement intervals with no correction. Raised to a simultaneous 95\% level,
that is 99.79\% each (a correction written after E2 and before this computation), they leave 12~settled changes instead of 14 (\cref{tab:e2-bonferroni}).
Two pairs lose their status: PPO JAX against CrossQ on \task{Hopper}, PPO PyTorch against SAC on
\task{Walker2d}. The verdict of E2 survives the correction.

\gentable{e2}

\gentable{e2-bonferroni}

\paragraph{What the preregistered count counts}
The number 14 counts triples, not pairs that reverse. For each task $T$ where a program $X$ leads by IQM,
the script counts every other task $U$ and every program $Y$ that beats $X$ in a settled way on $U$
(\file{sim/analysis/e2.py}, lines 34 to 37). It never asks whether $X$ beats $Y$ on $T$. \enquote{SAC leads
on \task{HalfCheetah} and loses to CrossQ on \task{Walker2d}} counts there as one settled change, although SAC
and CrossQ are not separated on \task{HalfCheetah}. Written after the results and before this
computation, a strict count fixes this defect. A pair $(A, B)$ reverses when $A$ beats $B$ in a settled way on one task
and $B$ beats $A$ in a settled way on another. \Cref{tab:e2-strict} contrasts the two counts. With ten
random seeds, three pairs reverse:
\begin{itemize}
  \item PPO PyTorch beats SAC on \task{Hopper}, \task{Walker2d} and \task{Swimmer}; SAC beats it on
    \task{HalfCheetah};
  \item PPO JAX beats SAC on \task{Hopper} and \task{Swimmer}; SAC beats it on \task{HalfCheetah};
  \item PPO JAX beats CrossQ on \task{Hopper}; CrossQ beats it on \task{HalfCheetah} and \task{Walker2d}.
\end{itemize}
The first two survive Bonferroni's correction. The preregistered verdict holds under the strict reading, on
three pairs instead of fourteen changes: the order of two programs changes in a settled way from one task to
another.

\gentable{e2-strict}

\paragraph{The learning curves}
\Cref{fig:courbes-apprentissage} shows the evaluation curves of the sixteen program-task pairs. It recalls
what a final ranking hides: the curves of SAC and CrossQ stop at \num{200000}~steps, those of PPO continue to
one million. The equal-budget reading then compares the four programs at the same point of these curves.

The ranking forms late. An exploratory analysis, written after the results and before this computation,
compares with twenty random seeds the ranking of the four programs at
step $k$ with the final ranking through Kendall's $\tau$, which equals 1 for two identical rankings and $-1$
for two reversed ones. Its median over the four tasks is 0.17 at \num{50000}~steps, 0 at \num{100000}, 0.17
at \num{150000} and 0.67 at \num{200000}. The program that leads at the end already leads at
\num{200000}~steps on two tasks out of four. At the level of a single run, the Spearman correlation (the correlation of
ranks) between the score at \num{50000}~steps and the final score has a median of $-0.07$ for PPO and 0.19 for SAC and
CrossQ: a run that looks promising early seldom stays ahead. These comparisons read PPO early in its budget of
one million steps.

\fig{courbes-apprentissage}{Evaluation curves of the four programs (PPO PyTorch, PPO JAX, SAC, CrossQ) on the
four tasks (one task per column). The horizontal axis gives the training steps, one point every
\num{50000}~steps: up to one million for PPO, \num{200000} for SAC and CrossQ. The vertical axis gives the
mean return of ten evaluation episodes played without randomness. The solid line is the interquartile mean of
the ten random seeds, the band its 95\% bootstrap interval, the faint lines the individual runs.}

\paragraph{Twenty random seeds}
With random seeds 0 to 19, the leading program still changes with the task: CrossQ leads on
\task{HalfCheetah} and \task{Walker2d}, PPO JAX on \task{Hopper} and \task{Swimmer}. There remain 10~settled
changes, 6 with Bonferroni's correction. Under the strict count, four pairs reverse, two with Bonferroni: the
two pairs of SAC with the PPOs, as before, plus CrossQ against PPO JAX and the pair of the two PPOs
(\cref{tab:e2-strict}). Three results diverge from those at ten random seeds.
\begin{itemize}
  \item On \task{HalfCheetah}, SAC yields the lead to CrossQ: its IQM goes from \num{5191} to \num{4214}, that
    of CrossQ from \num{4801} to \num{4694}. The difference between the two is settled neither at ten nor at twenty
    random seeds.
  \item On \task{Hopper}, CrossQ rises from third to second place (\num{2302}, then \num{2792}). The lead of PPO
    JAX over CrossQ is no longer settled.
  \item The leading program now beats all the others in a settled way on two tasks, \task{Walker2d} and
    \task{Swimmer}.
\end{itemize}

\paragraph{At equal budget}
PPO receives five times more steps than SAC and CrossQ. Read at \num{200000}~steps on its curve, a
reading written before any look at these scores, PPO no longer leads on any task.
\begin{itemize}
  \item With ten random seeds, SAC leads on \task{HalfCheetah}, CrossQ on the other three tasks. The
    preregistered count leaves 3~settled changes: SAC, leading on \task{HalfCheetah}, is beaten in a settled
    way by CrossQ on \task{Walker2d} and \task{Swimmer} and by PPO JAX on \task{Swimmer}. On
    \task{Hopper}, no pair is settled. Under the strict count, two pairs reverse, with or without Bonferroni:
    the pair of the two PPOs and that of PPO JAX against SAC.
  \item With twenty random seeds, CrossQ leads on all four tasks. The preregistered count finds no settled
    change at all, since no program beats the leading program. On \task{HalfCheetah}, its lead over SAC (\num{4694}
    against \num{4214}) is not settled either.
\end{itemize}
The leading program no longer changes, yet under the strict reading the order of two pairs still changes in a settled
way between tasks.
\begin{itemize}
  \item SAC beats PPO JAX on \task{HalfCheetah}, with a probability of improvement of 1.00 $[1.00, 1.00]$, as well as
    on \task{Walker2d}, with 0.99 $[0.97, 1.00]$. PPO JAX beats SAC on \task{Swimmer}, 0.91 $[0.80, 0.99]$.
    This reversal survives Bonferroni, with the interval $[0.73, 1.00]$ on \task{Swimmer}.
  \item PPO PyTorch beats PPO JAX on \task{HalfCheetah}, 0.72 $[0.55, 0.87]$. PPO JAX beats PPO PyTorch on
    \task{Swimmer}, 0.90 $[0.78, 0.99]$. This reversal does not survive Bonferroni.
\end{itemize}

\paragraph{What E2 establishes and its limits}
The preregistered verdict holds: with ten random seeds and at the protocol's budgets, the leading program
changes from task to task, in a settled way, even after Bonferroni's correction. The preregistered count
inflates the evidence. Three pairs really reverse, two with Bonferroni, instead of fourteen changes. Three
qualifications weigh on the verdict.
\begin{itemize}
  \item The leading program seldom beats all the others in a settled way.
  \item With twenty random seeds, the change of leader from one task to the next owes much to PPO's budget,
    five times larger. At equal budget, CrossQ leads everywhere. The order of the other programs still changes
    with the task at equal budget, for PPO JAX against SAC in a settled way even with Bonferroni. A paper that
    compares two methods at different budgets also compares these budgets.
  \item The ranking applies to the libraries' default values. Tuning per task might change it, in a direction
    we do not know. In one case, experiment E6 shows a published tuning doing worse than the default values.
\end{itemize}

\section{E3: two programs of the same algorithm part ways on one task}\label{sec:e3}

\paragraph{Question and plan}
Do two programs that implement the same algorithm give the same result? E3 sets PPO PyTorch against PPO JAX on
the four tasks, with the same random seeds, the same algorithm hyperparameters and one million steps
(\cref{sec:plan-e3}). The two programs differ in their numerical library and in how they initialise the
weights (\cref{sec:programmes}).

\paragraph{Ten random seeds: no winner, large differences}
\Cref{tab:e3} sets out the two IQMs, their ratio and the probability that a PPO JAX run beats a PPO PyTorch
run. On no task does either PPO win under the settled-winner rule. Yet the apparent differences are large. On
\task{Walker2d}, the IQM is \num{2158} with PyTorch and \num{1272} with JAX, 41\% lower. The probability that
JAX does better is 0.27, from 0.06 to 0.52: the interval still touches 0.5. The two IQM intervals overlap. On
the other three tasks, the probabilities of improvement stay close to 0.5: 0.36 on \task{HalfCheetah}, 0.54
on \task{Hopper}, 0.66 on \task{Swimmer}. The IQM intervals of E3 differ by a few points from those of E2
for the same runs, because each script consumes the random generator in its own order
(\cref{sec:reproductibilite-bootstrap}).

\gentable{e3}

\paragraph{Twenty random seeds: the two PPOs part ways}
Of all the results, this one moves most when random seeds are added. With twenty random seeds, two tasks have
a settled winner:
\begin{itemize}
  \item On \task{Walker2d}, PyTorch wins, \num{2204} against \num{1435}. The probability that JAX does better
    is 0.30, from 0.15 to 0.48.
  \item On \task{Swimmer}, JAX wins, 115 against 107. The probability is 0.79, from 0.63 to 0.91, with
    disjoint IQM intervals.
\end{itemize}
The differences resemble those at ten random seeds. The ten added random seeds narrow the intervals.

\paragraph{Four tasks, one correction}
The settled-winner rule judges each task alone, whereas E3 tests four tasks and the twenty-seed reading looks
at the same question a second time. Written after the results and before this computation, a complement adds a Mann-Whitney test
per task, corrected by Holm's procedure over the four tasks. It also gives the difference in score points through
the Hodges-Lehmann estimator, the median of the differences between a PPO JAX run and a PPO PyTorch run
(\cref{tab:e3-holm}). With ten random seeds, no task passes the correction. With twenty:
\begin{itemize}
  \item on \task{Swimmer}, PPO JAX wins even after correction ($p = 0.002$, Holm $p$ 0.009), by
    9~points $[3, 15]$ for IQMs of about 110: a settled effect, small in value;
  \item on \task{Walker2d}, the lead of PPO PyTorch does not pass the correction ($p = 0.029$, Holm $p$
    0.088). At the simultaneous level of 98.75\%, the probability that JAX does better is 0.30
    $[0.12, 0.54]$, an interval that contains 0.5. The difference is $-736$~points $[-\num{1267}, -116]$, a 95\%
    interval with no correction.
\end{itemize}
The test lacks power for a difference like the one on \task{Walker2d} (\cref{tab:e3-puissance}). At the corrected
level $0.05/4$, it detects that difference in 14\% of the draws with ten random seeds, 37\% with twenty, 59\% with
thirty, 76\% with forty and 92\% with sixty. At 80\% power, the smallest detectable probability of improvement
is 0.94 with ten random seeds, 0.81 with twenty and 0.72 with forty.

\gentable{e3-holm}

\gentable{e3-puissance}

\paragraph{One difference survives the correction}
With ten random seeds, we cannot say that the two PPOs differ. With twenty, the preregistered rule
separates them on two tasks, in opposite directions. After correction for the four tasks, one difference
remains settled: PPO JAX wins on \task{Swimmer}, by 9~points. On \task{Walker2d}, the lead of PyTorch remains
likely without being settled. The test has only about one chance in three of seeing such a difference with
twenty random seeds. Flaw F14, the implementation that changes the result, therefore shows up again in 2026 on
one task, between two libraries by the same author. These two libraries share their algorithm hyperparameters but
not their default initialisation. The option \texttt{ortho\_init} is true in Stable-Baselines3 and false in
SBX, which then keeps Flax's initialisation for every layer. When switched on, the option does not do the same
thing in the two libraries either. In SBX, it initialises only the action output layer, at scale 0.01. In Stable-Baselines3, it also initialises the hidden layers, with gain $\sqrt2$, then the
value output, with gain 1 (\cref{sec:programmes}). E3 does not identify
which code difference causes the difference in scores: the weight initialisation, the random generators or the rounding of the
numerical library. An experiment written in advance adds twenty random seeds on \task{Walker2d} and
\task{Swimmer}, then swaps the initialisation of the two programs. It is still running.

\section{E4: keeping the best three out of ten inflates the mean by a third}\label{sec:e4}

\paragraph{Question and plan}
What happens to a result when only the best three of ten runs are published, a practice
\textcite{henderson2018matters} pointed out? E4 applies this selection to each of the sixteen program-task
pairs of E2, then counts the comparisons, out of 24, that change direction (\cref{sec:plan-e4}).

\paragraph{A mean always inflated}
\Cref{tab:e4} and \cref{fig:meilleurs-sur-dix} give the inflation for each pair. All sixteen inflations are
positive, by construction. The mean inflates by 32\% at the median, from 6\% to 84\% depending on the program
and the task. The largest hit PPO JAX on \task{Walker2d} (84\%), SAC on \task{Hopper} (78\%) and CrossQ on
\task{HalfCheetah} (74\%). The standard deviation of the best three has a median of 35\% of that of all ten:
the selection inflates the mean while it narrows the apparent spread, which makes the result look more
certain.

\fig{meilleurs-sur-dix}{Rise of the mean when only the best three of ten runs are published, for each program
(one colour) and each task (one row), as a percentage of the mean of all ten. The dotted line marks the
median, +32\%. All sixteen rises are positive: sorting always serves the one who sorts. Data: ten random seeds
per pair, final evaluation scores.}

\paragraph{Three comparisons reversed}
Out of 24~comparisons, 3 change direction when each program keeps only its best three runs: PPO PyTorch
against PPO JAX on \task{HalfCheetah} and on \task{Hopper}, CrossQ against PPO PyTorch on \task{Swimmer}. Two
of the three set the two PPOs against each other. When the difference is small, keeping one's best runs
changes the winner.

\gentable{e4}

\paragraph{Replication on random seeds 10 to 19}
Repeated identically on ten other random seeds, the selection inflates the mean by 43\% at the median, from
7\% to 85\%. It again reverses 3~comparisons out of 24: CrossQ against PPO PyTorch on \task{Hopper}, PPO
PyTorch against PPO JAX on \task{Hopper}, CrossQ against PPO PyTorch on \task{Swimmer}. Two of the three
reversals are the same as on random seeds 0 to 9.

\paragraph{What E4 establishes}
Selecting the best runs always inflates the mean, by about a third at the median. It reverses close
comparisons. The result replicates on independent random seeds. E4 measures the effect of a practice. How many
papers follow the practice is a question for the survey (criterion X3, \cref{chap:04-sondage-resultats}).

\section{E5: five random seeds are not enough for a small difference}\label{sec:e5}

\paragraph{Question and plan}
How many random seeds does it take to recover the correct order of two programs? E5 draws $n$ runs per program
with replacement from those observed, \num{10000} times, then counts the share of draws that recover the order
of the means of all the runs (\cref{eq:e5}, \cref{sec:plan-e5}).

\paragraph{Ten random seeds}
Over the 24 comparisons, the median share of draws that recover the correct order is 94\% with three random
seeds, 97\% with five and 99\% with ten (\cref{tab:e5}). This reassuring figure comes from the large
differences between our programs. Seven comparisons are close, with means that differ by less than 20\%. For them, the
median falls to 65\% with three random seeds, 71\% with five, 77\% with ten and 85\% with twenty. The hardest
comparison, CrossQ against SAC on \task{HalfCheetah}, recovers the correct order in only 60\% of the draws with
five random seeds.

\gentable{e5}

\paragraph{Twenty random seeds}
The correct order becomes that of the means of twenty runs. The draws are made among the twenty (\cref{fig:e5-graines}, right panel). Over the
24 comparisons, the median barely moves: 93\% with three random seeds, 97\% with five. The close comparisons
now number eight. They are harder than ten random seeds suggested: the median is 59\%, 62\%, 67\% and 74\%
with three, five, ten and twenty random seeds. CrossQ and PPO PyTorch on \task{Swimmer} have means that differ
by less than 1\%: even twenty random seeds order them in only 53\% of the draws, about as often as a coin
toss.

\fig{e5-graines}{Share of the \num{10000} draws that recover the correct order of two programs, as a function
of the number of random seeds drawn for each (3, 5, 10, 20). The correct order is that of the means of all
the observed runs. One curve per comparison between two programs on one task, 24 in all. In red, the close
comparisons, whose means differ by less than 20\%. Two panels: the correct order and the draws come from ten
random seeds, then from twenty. With twenty random seeds as the reference, five recover the correct order of
the close comparisons in only 62\% of the draws at the median.}

\paragraph{Checked out of sample}
E5 takes as the correct order the sign of the means of the same sample it draws from, with replacement. Its
share therefore measures the stability of the observed sign, not the probability of recovering the true
order. For the same reason, the twenty-seed column of the ten-seed panel draws twenty runs from ten: it is an
extrapolation. A check written after the results and before the computation removes this circularity. It judges the order
of random seeds 0 to 9 on random seeds 10 to 19, then the reverse (\cref{tab:e5-hors-echantillon}). The sign
of the difference between two programs is found again in 20 comparisons out of 24, including 5 of the 7 close
comparisons. Four signs flip: PPO PyTorch against PPO JAX on \task{HalfCheetah} and CrossQ against PPO PyTorch
on \task{Swimmer}, two close comparisons, as well as CrossQ against each of the two PPOs on \task{Hopper}. The lead
of PPO JAX over CrossQ on \task{Hopper}, settled at ten random seeds with the interval $[0.52, 0.94]$, does not
recur on the next ten random seeds. Drawn without replacement from random seeds 0 to 9
and judged on random seeds 10 to 19, $n$ runs recover the order of the close comparisons in 60\% of the draws
on average with three random seeds, 65\% with five and 71\% with ten. For the same comparisons, E5 gave 65\%,
71\% and 77\% at the median. Without circularity, the shares therefore stay close to those of E5, slightly
lower in this direction. In the other direction they are not lower: 69\% against 65\% with three random
seeds, 71\% against 71\% with five. With ten, all ten runs of a half are drawn, so each comparison is
recovered in every draw or in none. The share is then 5 of 7 in both directions. E5 is a little optimistic in
one direction only.

\gentable{e5-hors-echantillon}

\paragraph{Close comparisons need more than five random seeds}
Five random seeds, the median value the survey finds in the papers of 2026 (\cref{chap:04-sondage-resultats}),
suffice for large differences. They do not suffice for differences under 20\%: with twenty random seeds as the
reference, such a difference is read backwards more than one time in three. The out-of-sample check confirms this order of
magnitude. E5 lets the observed runs stand in for the true distribution. With ten or twenty runs, this
distribution is poorly known. For the closest comparisons, so is the correct order.

\section{E6: a published recipe has aged with its libraries}\label{sec:e6}

\paragraph{Question and plan}
E6 was meant to illustrate flaw F7: the lead a new method credits to itself when it alone is tuned. On
\task{Hopper-v5}, E6 sets ten PPO PyTorch runs with the default values of Stable-Baselines3 against ten runs
with the settings published by RL Baselines3 Zoo, the collection of settings that serves as the reference
for this library, called the Zoo below (\cref{sec:plan-e6}). \Cref{tab:e6} gathers every group.

\gentable{e6}

\paragraph{First result: the published settings would do three times worse}
The default values reach an IQM of \num{2877} $[\num{2268}, \num{3369}]$. The tuned recipe reaches
\num{1003} $[856, \num{1146}]$. In means, \num{2776} against 994, with a plateau near \num{1000} reached as
early as \num{50000}~steps. The published settings would do nearly three times worse than the default values.

\paragraph{An error on our side that explains nothing}
The result was too good to believe without a check. We found an error in our code. The Zoo's reward
normalisation uses the algorithm's discount factor, 0.999. Our first pass left it at its default value, 0.99.
The diagnosis, decided after this first result, reruns the corrected recipe on
\task{Hopper-v5} and on \task{Hopper-v4}, plus the default values on v4, ten random seeds each.
\begin{itemize}
  \item The corrected recipe obtains \num{1012} $[902, \num{1135}]$ on v5 and \num{984} $[773, \num{1002}]$ on
    v4, the version the Zoo was tuned for.
  \item The default values obtain \num{2423} $[\num{1657}, \num{3058}]$ on v4.
  \item The probability that a corrected run beats a run with the error is 0.47, from 0.21 to 0.73: the error
    changes nothing detectable.
  \item On v4, a corrected Zoo run beats a default run with a probability of only 0.07, from 0 to 0.24. The
    difference survives Holm's correction ($p=0.001$). On v5, the same probability is 0.01.
\end{itemize}
Our error was real, yet it does not explain the difference. Moving from v5 to v4 does not explain it either:
on v4 as on v5, the recipe plateaus near \num{1000}.

\paragraph{The official script plateaus too}
Two explanations remained: the published recipe really does worse than the default values with today's
libraries, or our way of applying it still differs from the Zoo's. A control written before any
official training decides between them with the Zoo's own script. The ten policies trained by
\texttt{rl\_zoo3.train} on \task{Hopper-v4}, evaluated with our protocol, obtain \num{995}
$[821, \num{1006}]$. This matches our corrected recipe, \num{984} $[777, \num{1002}]$ in this analysis: the
probability that an official policy beats ours is 0.58, from 0.32 to 0.82. The default values obtain
\num{2423} $[\num{1687}, \num{3058}]$. An official policy beats a default policy with a probability of only
0.09, from 0 to 0.30. The difference survives Holm's correction ($p=0.001$). The intervals of \num{984} and
\num{2423} differ by a few points between the two analyses, because each draws its own resamples.

\paragraph{A robot that stands without moving}
On \task{Hopper}, the reward of one step is
\begin{equation}
  r_t = v_{x,t} + 1 - 10^{-3}\,\lVert a_t\rVert^2,
\end{equation}
where $v_{x,t}$ is the forward velocity in metres per second, the term 1 is the bonus paid while the robot
stays upright and $a_t$ is the action (Gymnasium's default values). A step lasts 0.008~seconds, so
\num{1000}~steps take eight seconds. A return close to \num{1000} over \num{1000}~steps is that of a robot
that stands upright without moving forward. Two controls, each written before its computation, checked
this.
\begin{itemize}
  \item The official training logs. Over the last 200 episodes of each random seed, the
    episodes that reach \num{1000}~steps earn a median of 1.002 per step, from 0.978 to 1.013 depending on the
    random seed: the bonus alone, with no forward progress.
  \item The final policies replayed with our protocol (\cref{tab:e6-vitesse},
    \cref{fig:e6-vitesse}). The thirty replayed policies, official, corrected Zoo and default, give back the
    returns already published. The official agents
    move forward at 0.00~m/s at the median, from $-0.45$ to 0.03~m/s. Eight random seeds out of ten hold all
    ten episodes up to the \num{1000}th~step. Random seed 0 falls within about thirty steps. The agents of the
    corrected Zoo recipe behave the same way. Our default agents run at 2.45 to 2.89~m/s.
\end{itemize}
The evaluation videos show it (\cref{fig:e6-images}). The official agent of random seed 7 stays upright for
\num{1000}~steps and ends 2~centimetres behind its starting point. The default agent of random seed 5 runs
19~metres in 862~steps, then falls. The Zoo's evaluation script, \texttt{rl\_zoo3.enjoy}, also gives
\num{996.55} over five episodes of \num{1000}~steps for random seed 7.

\gentable{e6-vitesse}

\fig{e6-vitesse}{Forward speed of the PPO PyTorch agents on \task{Hopper-v4}, one dot per random seed,
replayed with our protocol (ten episodes without randomness). The agents trained by the official script of
RL Baselines3 Zoo and those of its recipe applied by our code stand upright without moving forward
(median speed 0.00~m/s, return near \num{1000} over \num{1000}~steps). The agents with the default values run
at 2.45 to 2.89~m/s. The agent the Zoo published in 2022, with the same settings and the libraries of the
time, moves forward at 1.44~m/s. On the right, the median return of each group.}

\fig{e6-images}{Four frames (steps 0, 240, 480 and 720) of two evaluation videos on \task{Hopper-v4}, with a
camera that follows the robot. Top, the official Zoo agent, random seed 7: it stays upright on the same square
of the checkerboard for \num{1000}~steps, for a return of 997. Bottom, the agent with the default values,
random seed 5: it runs 19~metres in 862~steps before falling.}

\paragraph{The recipe worked in 2022}
One last control, whose reading was fixed in advance, replays the agent the Zoo itself published for this
task, \texttt{sb3/ppo-Hopper-v3} on Hugging Face (revision \texttt{b29cb57\dots}). Its \file{config.yml} file gives our hyperparameters. It was trained in 2022
with Stable-Baselines3~1.5.1a8, Gym~0.21 and PyTorch~1.11. Its model card reports \num{2410}~$\pm$~10 on
\task{Hopper-v3}. Replayed with our protocol on \task{Hopper-v4}, it obtains \num{2439} over ten episodes and
\num{2438}~$\pm$~4 over a hundred, all of \num{1000}~steps, at 1.44~m/s.
The reading fixed in advance applies: the recipe worked with its original libraries and has aged with
them. This control rests on a single published agent, perhaps the best of a tuning run, whereas our ten
official random seeds all plateau near \num{1000}. A file serialized by Python can run code when it is
loaded. Every object in the file was therefore inspected before any loading. The file contains only tensors,
NumPy arrays, a Gym space and Stable-Baselines3 classes.

In score points, the Hodges-Lehmann estimator, added before this computation, gives
the typical difference between an official policy and a default policy on v4: $-\num{1464}$~points
$[-\num{2342}, -606]$. Between an official policy and a policy of the corrected Zoo recipe, it is $+5$~points
$[-18, 54]$. The difference from the default values is large. The difference from our application of the
recipe is zero to the precision of the measurement.

\paragraph{A published tuning that has aged}
Under the reading rule fixed in advance, our application of the recipe was not at fault. With
Stable-Baselines3~2.9.2, Gymnasium~1.1.1 and MuJoCo~4.0.0, the recipe tuned and published by the Zoo
does worse than the default values on \task{Hopper}, on version 4 as on version 5. It produces a robot that
stands still. Applied in 2022 with Stable-Baselines3~1.5, the same recipe had produced an agent that runs, at
\num{2439} under our protocol. E6 does not show a baseline its authors tuned badly. It shows a published tuning
that has aged with the software around it, to the point of doing worse today than the default values of the
same library. E6 also recalls that a result that looks too good deserves a check that goes as far as the
official script.

\paragraph{Limits}
We do not know which library change made the recipe age. We tried only one task and one algorithm of the
Zoo. The 2022 control rests on a single published agent. A rerun of the recipe with the discount factor
0.99 instead of 0.999 is written in advance and still running.

\section{E7: in a nearby physics, the leader on Hopper ends last}\label{sec:e7}

\paragraph{Question and plan}
Pineau asked that training be kept apart from evaluation. Our agents are already evaluated on random seeds
never seen in training. E7 goes further: each final policy also plays in a modified physics, with masses made
20\% heavier and sliding friction reduced by 20\%. For each run, we compute the share of score kept, the ratio
of the score in the modified physics to the score in the training physics, then its IQM over the ten random
seeds with its interval (\cref{sec:plan-e7}).

\paragraph{Twelve pairs out of sixteen lose}
\Cref{tab:e7} and \cref{fig:e7-carte} compare the share kept from one program-task pair to the next. Twelve pairs out of
sixteen lose a settled share of their score: their interval stays below 1. Three change by no
detectable amount. One gains: SAC on \task{Swimmer}, at 1.07 $[1.02, 1.12]$. The loss depends on the program
and the task.
\begin{itemize}
  \item On \task{Hopper}, PPO keeps only 22\% of its score in its PyTorch version and 18\% in its JAX version,
    while SAC keeps 64\% and CrossQ 65\%.
  \item On \task{Walker2d}, PPO PyTorch keeps 34\%, CrossQ 46\%, PPO JAX 52\%, SAC 84\%.
  \item On \task{HalfCheetah}, the losses are smaller: from 85\% kept for PPO PyTorch to 92\% for PPO JAX.
  \item On \task{Swimmer}, no program loses.
\end{itemize}

\fig{e7-carte}{Share of score kept in a modified physics (masses multiplied by 1.2, sliding friction by 0.8),
for each program and each task. Each cell gives the interquartile mean, over ten random seeds, of the ratio of
the score in the modified physics to the score in the training physics, both measured on the same evaluation
random seeds. Tasks are in columns, programs in rows. A cell marked as a loss has a 95\% interval entirely
below 1, a cell marked as a gain an interval above 1.}

\gentable{e7}

\paragraph{The ranking upended on Hopper}
A descriptive complement added after the results, with no test, ranks the programs in each physics by the IQM of
their score. On \task{Hopper}, the order reverses. PPO JAX, first in the training physics with \num{3064},
ends last in the modified physics with 525. CrossQ, third with \num{2302}, takes the lead with \num{1302},
ahead of SAC (751) and PPO PyTorch (609). On \task{Walker2d}, SAC moves from last place to second. On
\task{HalfCheetah} and \task{Swimmer}, the order does not change. \Cref{fig:physique-modifiee} shows the
reversal on \task{Hopper}.

\fig{physique-modifiee}{\task{Hopper}, four programs, ten random seeds each. On the left, the score in the
training physics. On the right, the score with masses made 20\% heavier and friction reduced by 20\%. Each
line links the score and rank of one program in the two physics. PPO written in JAX goes from first to last
place, CrossQ from third to first. The scores are interquartile means of the final evaluation scores. The
ranks follow from these means. \Cref{tab:e7-test} tests the pairs one by one.}

\paragraph{The reversal, tested}
A ranking by IQM does not say which pairs really change order. A complementary test judges each pair on
\task{Hopper} through its probability of improvement, in each physics, following a plan written before the computation
(\cref{tab:e7-test}).
A pair reverses in a settled way when it is settled in one direction in the training physics and in the other
direction in the modified physics.
\begin{itemize}
  \item In the modified physics, CrossQ beats each of the three other programs in a settled way, at ten as at
    twenty random seeds. The probability that a PPO JAX run beats a CrossQ run there is 0.07 $[0.00, 0.24]$
    with ten random seeds (Mann-Whitney test, $p = 0.001$) and 0.04 $[0.00, 0.12]$ with twenty. The
    probability that it beats SAC is 0.11 $[0.00, 0.28]$, then 0.12 $[0.03, 0.24]$.
  \item With ten random seeds, three pairs reverse in a settled way: PPO JAX against CrossQ (0.77
    $[0.52, 0.94]$ in the training physics, 0.07 in the modified physics), PPO JAX against SAC (0.89, then
    0.11) and PPO PyTorch against SAC (0.89, then 0.25 $[0.07, 0.49]$).
  \item With twenty random seeds, two pairs reverse: the two PPOs against SAC. The lead of PPO JAX over CrossQ
    in the training physics is no longer settled, at 0.56 $[0.37, 0.73]$, as the out-of-sample check of E5
    suggested.
  \item The \enquote{last place} of PPO JAX is not settled against PPO PyTorch at ten random seeds, at 0.65
    $[0.41, 0.88]$ for PyTorch. It is settled, narrowly, at twenty, at 0.68 $[0.51, 0.84]$.
  \item On \task{Walker2d}, PPO PyTorch against SAC also reverses in a settled way, at both sizes.
    \task{HalfCheetah} and \task{Swimmer} have no reversal.
\end{itemize}
This test is a complement decided after the results.

\gentable{e7-test}

\paragraph{Twenty random seeds}
The verdicts stay the same for 14~pairs out of 16. Two change. SAC on \task{Walker2d} no longer loses by a
detectable amount: 0.91, from 0.79 to 1.01. CrossQ on \task{Swimmer} gains a little: 1.03, from 1.01 to 1.04.
That makes 11~losses and 2~gains. The descriptive ranking still reverses on \task{Hopper}, the only task where
the leading program changes.

\paragraph{What E7 establishes and its limits}
On \task{Hopper}, the program that leads in the training physics ends far behind in a neighbouring physics:
CrossQ beats each of the three other programs in a settled way in the modified physics, at ten and twenty
random seeds, including PPO JAX, which led in the training physics. A ranking obtained in the training physics
says little about robustness. E7 tests only one modification, which changes two quantities at once.
\Cref{sec:dose} separates the two quantities and varies the size of the change. The share kept is computed run by
run, so it does not read as the ratio of the two IQMs of the ranking.

\section{The physics by degrees: the loss grows with the size of the change}\label{sec:dose}

\paragraph{Question and plan}
Does the reversal of E7 hinge on one modification chosen in advance? An addition written before any computation
re-evaluates each final policy while changing one quantity at a time: masses from 30\% lighter to 30\% heavier, friction from
40\% weaker to 20\% stronger. A pair loses when the interval of the IQM of its ratio to the score in the training
physics stays below 1, as in E7 (\cref{sec:plan-e7}).

\paragraph{The control}
At 1.0 times, the re-evaluations recover the final score written by the training script: the largest relative difference
is zero, for the 160 policies at ten random seeds as for the 320 at twenty.

\paragraph{The pairs that lose}
\Cref{tab:dose}, \cref{fig:intensites} and \cref{fig:dose-carte} count the pairs that lose, out of 16, for each
physics. With ten random seeds:
\begin{itemize}
  \item masses multiplied by 0.7, 0.8, 0.9: 8, 6 and 3~pairs lose;
  \item masses multiplied by 1.1, 1.2, 1.3: 9, 12 and 11~pairs lose;
  \item friction multiplied by 0.6, 0.7, 0.8, 0.9: 7, 6, 4 and 3~pairs lose;
  \item friction multiplied by 1.1 and 1.2: 1 and 9~pairs lose.
\end{itemize}
The losses trace a V around the training physics, steeper on the side of heavier masses. Making the masses
20\% heavier alone makes 12~pairs lose, as many as the combined modification of E7. Reducing friction by 20\%
makes only 4 lose. In E7, the loss seems to come mostly from the masses, although neither experiment tests
this decomposition. A few pairs gain, mostly when friction drops: 3 gains at 0.6 times and at 0.7 times.

\gentable{dose}

\fig{intensites}{Four programs on four tasks, ten random seeds each, or 16 program-task pairs. Each point
counts the pairs that lose a detectable share of their score when one quantity of the physics changes: the
masses (from 30\% lighter to 30\% heavier) or the sliding friction (from 40\% weaker to 20\% stronger). A
pair loses when the 95\% interval of its score, divided by its score in the training physics, is entirely below 1.
At 0\%, the training physics serves as the reference: no loss can occur there.}

\fig{dose-carte}{Share of score kept in twelve modified physics, with ten random seeds, then twenty, shown next
to each other. One row per program and task (16), one column per physics: masses from 70\% to 130\% of their
training value, then sliding friction from 60\% to 120\%, one quantity changed at a time. A cell is coloured
only if the 95\% interval of the interquartile mean of the ratio to the score in the training physics avoids 1: red for a
loss, darker when more than half the score is lost, green for a gain. Under each column, the number of pairs
that lose, out of 16.}

\paragraph{The leading program on Hopper, physics by physics}
On \task{Hopper}, PPO JAX, first in the training physics, is no longer first in 10 of the 12 modified physics.
It keeps the lead only with friction at 0.9 and 1.1 times. The lead passes to PPO PyTorch or to CrossQ
depending on the quantity and the factor. On the other three tasks, the leading program changes little: SAC in
all twelve physics of \task{HalfCheetah}, CrossQ in all twelve of \task{Walker2d}, PPO JAX in ten of the twelve
of \task{Swimmer}.

\paragraph{Twenty random seeds}
The V keeps its shape. With masses at 0.9, 1.1, 1.2 and 1.3 times, 4, 11, 12 and 12~pairs lose. With friction
at 1.2 times, 10. On \task{Hopper}, PPO JAX again loses the lead in 10 physics out of 12. On
\task{HalfCheetah}, CrossQ takes the lead in 9 of the 12 physics, as it does in the training physics at twenty
random seeds.

\paragraph{Hopper's pairs, tested at each intensity}
The same test also judges the pairs on \task{Hopper} at each intensity, with twenty random seeds. CrossQ
beats each of the three other programs in a
settled way when friction drops to 0.6, 0.7 or 0.8 times and when the masses rise to 1.3 times. With masses at
1.1 and 1.2 times, or friction at 0.9 times, CrossQ leads and beats PPO JAX in a settled way without beating
all the others. With masses at 0.7 times, PPO PyTorch beats the other three. The reversal on \task{Hopper}
comes mostly from reduced friction and heavier masses. Reduced friction makes few of the sixteen pairs lose,
yet it reorders \task{Hopper}.

\paragraph{What the physics by degrees establishes}
The reversal on \task{Hopper} does not hinge on one chosen modification. It appears in almost every
neighbouring physics, at ten as at twenty random seeds. The rankings by IQM have no test on the other three
tasks. On \task{Hopper}, the pairwise test settles the lead of CrossQ in four of the twelve modified physics.
The modifications are small and uniform: each mass or each friction coefficient changes by the same factor.
They cannot replace a real test of generalisation, on other bodies or other tasks.

\section{B3: the chosen measure changes the winner of close comparisons only}\label{sec:b3}

\paragraph{Question and plan}
Can a paper that does not define its measure choose its winner? B3 ranks the four programs by three measures
computed on the same curves: the final score, the best score reached along the curve and the mean of the curve
(\cref{sec:plan-b3}).

\paragraph{Ten random seeds: eight reversals, none settled}
Eight pairs of programs change order from one measure to another (\cref{tab:b3}). On two tasks, the leading
program changes with the measure.
\begin{itemize}
  \item On \task{HalfCheetah}, the mean of the curve puts CrossQ ahead of SAC, which leads by the final score
    and by the best score.
  \item On \task{Hopper}, the best score puts PPO PyTorch ahead of PPO JAX, which leads by the other two
    measures.
\end{itemize}
The other reversals involve CrossQ and PPO PyTorch on \task{Hopper} and on \task{Swimmer}. None is settled:
they set close programs against each other, which ten random seeds do not separate under any measure. The
mean of the curve covers one million steps for PPO and \num{200000} for SAC and CrossQ. It favours a program
that learns fast early on.

\gentable{b3}

\paragraph{Twenty random seeds and equal budget}
\begin{itemize}
  \item At unequal budgets with twenty random seeds, 8~pairs still change order, with no settled reversal. The
    leading program changes with the measure only on \task{Hopper}.
  \item At equal budget with ten random seeds, 4~reversals, none settled. The leading program changes with the
    measure only on \task{HalfCheetah}, between SAC and CrossQ.
  \item At equal budget with twenty random seeds, 6~reversals, of which 1 is settled: on \task{Swimmer}, PPO
    PyTorch and SAC change order between the best score reached and the mean of the curve. The leading
    program changes with the measure on no task.
\end{itemize}

\paragraph{What the best point adds}
The exploratory analysis of the curves also measures, run by run, what reporting the best point of the
curve instead of the final score gains. For PPO, the final score is below the best point in 55\% to 100\% of
the runs depending on the task. Reporting the best point then inflates the IQM by up to 48\% (PPO JAX on
\task{Walker2d}, ten random seeds), 37\% for the same pair at twenty and 33\% for PPO PyTorch on
\task{Hopper} at twenty. For SAC and CrossQ, evaluated at only four points, the inflation stays under 10\%,
except for CrossQ on \task{Hopper}: 41\% at ten random seeds, 17\% at twenty. A paper that reports the best
point without saying so inflates its long-curve programs first.

\paragraph{What B3 establishes}
Between two close programs, the chosen measure sometimes decides the apparent winner. Between clearly
separated programs, the measure changes nothing. Only one reversal is settled across the four analyses.
It involves two programs far from the leader. A paper that does not state its measure leaves the winner of a
close comparison an open question.

\section{Ten or twenty random seeds, equal budget: what holds and what changes}\label{sec:consolidation}

The results with ten random seeds remain the main results, as the protocol announced. The additions written
before any new training (twenty random seeds, equal budget, replication) confirm or qualify them.

\paragraph{What holds}
\begin{itemize}
  \item The leading program changes with the task at the protocol's budgets (E2), in a settled way, with or
    without Bonferroni's correction. Under the strict count, the order of four pairs reverses at twenty random
    seeds (three at ten), two of them with Bonferroni at both sizes.
  \item Selecting the best three runs out of ten always inflates the mean and reverses 3~comparisons out of
    24, on random seeds 10 to 19 as on random seeds 0 to 9 (E4).
  \item The modified physics upends the ranking on \task{Hopper}. The reversal of PPO JAX and PPO PyTorch
    against SAC remains settled at twenty random seeds. The verdicts of E7 stay the same for 14~pairs out of
    16. The physics by degrees keeps its V shape.
  \item The reversals of B3 at unequal budgets still number 8, with no settled reversal.
\end{itemize}

\paragraph{What changes}
\begin{itemize}
  \item On \task{HalfCheetah}, CrossQ takes the lead from SAC, with no settled difference between them (E2).
  \item The two PPOs part ways on \task{Walker2d} and \task{Swimmer}, in opposite directions, under
    the preregistered rule (E3). After Holm's correction for the four tasks, only the difference on
    \task{Swimmer} remains settled.
  \item The close comparisons are harder to order: with five random seeds, the correct order is recovered in
    only 62\% of the draws at the median, instead of 71\% (E5).
  \item At equal budget with twenty random seeds, CrossQ leads on all four tasks and the preregistered count
    finds no settled change (E2). Under the strict count, the order of two pairs still changes from one task
    to another, PPO JAX against SAC even with Bonferroni.
\end{itemize}
The last point is the qualification that weighs most, already stated in the summary of E2. The announced
verdict stands. With twenty random seeds, however, PPO's budget, five times larger, largely explains why the
leading program changes from one task to the next. The ranking of the programs that follow it still depends on
the task.

\section{What these experiments leave open}\label{sec:limites-experiences}

The limits specific to each experiment are stated in its section. A few limits apply to all of them.
\begin{itemize}
  \item All runs were made on one machine, a Mac with an Apple M2 Max chip, on the CPU. A rerun elsewhere
    would give other curves for the same random seeds.
  \item The programs keep the default values of their libraries. The rankings apply to these settings. They
    say nothing of the algorithms at their best.
  \item Four MuJoCo tasks, all with planar or nearly planar bodies, do not represent the reinforcement learning
    of 2026. The language-model part (\cref{chap:07-modeles-de-langue}) covers another share of it, at small
    scale.
  \item Ten evaluation episodes, played without randomness from ten fixed starting positions, measure a policy
    on a small part of the situations it could meet.
  \item Ten random seeds, then twenty, leave wide intervals for close comparisons. E5 puts a number on this.
  \item Several analyses were added after a result had been seen: the second statistic of E1, the diagnosis of
    E6, the descriptive ranking of E7. \Cref{sec:avant-apres} lists them. The text presents them as such.
  \item The analyses added last (deviations 21 to 31) were written after a technical review of the repository, which had
    already computed a draft version of several of these numbers outside the repository. Each deviation says
    so. The repository's scripts made the computation after the deviation was written. Their results are
    published whatever they turned out to be. On one point, the repository's computation differs from the
    draft: the strict count of E2 at equal budget with twenty random seeds gives one pair with Bonferroni, not
    two.
  \item Three experiments are written in advance and still running: the determinism of SAC and CrossQ,
    the Zoo recipe with the discount factor 0.99 and the cause of the difference between the two PPOs
    (deviations 32 to 34 of the simulation protocol, \cref{chap:09-ecarts}). This chapter reports none of
    their results.
\end{itemize}
